\documentclass[11pt]{article}
\usepackage[margin=0.85in]{geometry}
\usepackage{lmodern}
\usepackage{microtype}
\usepackage{enumitem}
\usepackage{listings}
\usepackage[T1]{fontenc}

\lstset{
  basicstyle=\ttfamily\small,
  columns=fullflexible,
  keepspaces=true,
  showstringspaces=false,
  breaklines=true,
  frame=single,
  xleftmargin=0.4em,
  xrightmargin=0.4em,
  aboveskip=0.55em,
  belowskip=0.55em
}

\pagestyle{empty}
\setlength{\parskip}{0.35em}
\setlength{\parindent}{0pt}
\setlist[enumerate,1]{leftmargin=1.4em, itemsep=0.75em, topsep=0.3em}
\setlist[enumerate,2]{leftmargin=1.5em, itemsep=0.35em, topsep=0.25em, label=\alph*.}
\setlist[itemize]{leftmargin=1.3em, itemsep=0.15em, topsep=0.2em}

\begin{document}

\begin{center}
{\LARGE\bfseries Midterm 1 --- Sample 2}\\[0.35em]
DATA 412/612
\end{center}

\vspace{0.25em}

Practice exam, questions only. Weeks 1--7: the Quiz 1 material,
plus \texttt{rowwise()} and \texttt{across()}, \texttt{case\_when()},
\texttt{readr}, exploratory plots, and tidyr.

Put your answers in an R Markdown document and keep the code visible
(\texttt{knitr::opts\_chunk\$set(echo = TRUE)} in the setup chunk).
You may use your notes and the course website.
You may not use Google, any other website, or an LLM.
Do not discuss the exam with anyone else.
Load the tidyverse.

\begin{enumerate}
\item \textbf{A function} (8 pts).
Write \texttt{finite\_mean(x)}. Drop every \texttt{NA} and every
non-finite value, then return the mean of what remains.
If nothing remains, return \texttt{NA}.
Use a pipe inside the function.
Try it on \texttt{c(1, NA, Inf, 3, NaN)}.

\item \textbf{\texttt{case\_when()} and \texttt{distinct()}} (8 pts).
In \texttt{mpg}, add \texttt{drive}: \texttt{"front"} when
\texttt{drv} is \texttt{"f"}, \texttt{"rear"} when it is \texttt{"r"},
and \texttt{"four"} when it is \texttt{"4"}.
Return the distinct pairs of \texttt{class} and \texttt{drive}, and
the number of cars in each pair.

\item \textbf{Import} (10 pts).
Read the text below with \texttt{read\_csv()}.
Keep \texttt{id} as character, including the leading zeros.
Parse \texttt{when} as a date. Leave \texttt{amount} numeric.
Show the problems \texttt{readr} reports, and print the tibble
after a successful parse of the rows that do have a real date.

\begin{lstlisting}
id,when,amount
007,2024-01-05,12
008,2024-02-11,NA
009,not-a-date,4
\end{lstlisting}

\item \textbf{Separate and widen} (10 pts).

\begin{enumerate}
\item Split \texttt{place} into \texttt{city} and \texttt{state}.
\item Turn \texttt{long} into a table with one row per city and one
column per quarter.
\end{enumerate}

\begin{lstlisting}
places <- tibble(place = c("Austin, TX", "Boise, ID", "Tampa, FL"))

long <- tibble(
  city = c("Austin", "Austin", "Boise", "Boise"),
  quarter = c("Q1", "Q2", "Q1", "Q2"),
  sales = c(10, 12, 7, 9)
)
\end{lstlisting}

\item \textbf{A faceted plot} (4 pts).
From \texttt{mpg}, plot \texttt{displ} against \texttt{hwy},
with one panel per \texttt{drv}. Color is not required.
Title the figure.
\end{enumerate}

\end{document}
